\section{Navier--Stokes via the Averaged-Equation Route}
\label{sec:goertzel-ns-state}

This section is a research-program lens on Goertzel's averaged-equation route to
3D Navier--Stokes global regularity. It is written to direct creative effort and
is candid about the current verdict: the route is, in Lean, \emph{provably
scaling-uncleared} --- it does not prove regularity, and the obstruction is the
supercritical-scaling wall. Every conjecture below is labelled. Lean identifiers
cite the live \texttt{FluidDynamics/NavierStokes} development.

\subsection{The creative bet}

The classical energy method controls $\int|u|^2\,dx$ and tries to close the
Beale--Kato--Majda (BKM) vorticity criterion from there. It fails at the
critical scaling: under the parabolic rescaling $(t,x)\mapsto(\lambda^2 t,
\lambda x)$ the $L^2$ energy is \emph{supercritical} in 3D, so it cannot dominate
the nonlinearity at small scales. The route's bet is to sidestep this by working
with \emph{averaged} (mean-field) equations: replace pointwise velocity norms
with smoother averaged quantities, prove the averaged momentum equation,
represent the Reynolds-stress closure, and transport the energy inequality and
the critical scaling through the average. If averaging buys enough regularity to
upgrade the diagnostic from supercritical $L^2$ to critical $L^3$, the BKM
criterion could close.

The Lean development encodes this bet as a structured \emph{promotion gate}
rather than as prose. A route is allowed to claim regularity only after clearing
three layers: the scaling verdicts (\verb|NavierScalingCheck|: criticalNorm,
supercriticalUpgrade, averagedEquationCompatibility), the challenge-world
interfaces (\verb|NavierScalingWorldInterface|, the ten required witnesses), and
the four averaged-equation obligations (\verb|NavierAveragedEquationObligation|:
averagedMomentumEquation, reynoldsStressClosure, energyInequalityCompatibility,
criticalScalingCompatibility). This is the generative contribution: the route is
turned into a finite, auditable list of obligations, so progress is measurable
and overclaiming is mechanically blocked.

\subsection{Why the route is provably scaling-uncleared}

The current energy/BKM route supplies only the supercritical $L^2$ diagnostic
and none of the upgrade machinery. Lean proves this exactly:
\begin{itemize}\itemsep1pt
  \item the diagnostic index is $2$ and $L^2$ is supercritical in 3D ---
    \verb|finiteEnergyVelocityLpIndex| $=2$,
    \verb|velocityLpScalingRegime 2 = .supercritical|, while the critical index
    is $3$ (\verb|criticalVelocityLpIndex|);
  \item the critical-norm gate is blocked and the supercritical-upgrade gate is
    blocked --- \verb|currentNavierEnergyBKMScalingStatus| with
    \verb|criticalNorm = .blocked|,
    \verb|supercriticalUpgrade = .blocked|,
    \verb|averagedEquationCompatibility = .notRepresented|;
  \item all four averaged-equation obligations are missing ---
    \verb|currentNavierEnergyBKMAveragedEquationCompatibilityStatus| has
    \verb|represented = []|;
  \item the three-layer gate does not clear ---
    \verb|currentNavierEnergyBKMRoute_regularityPromotionGate_refuted|,
    \verb|NavierRegularityPromotionGate.ClearsAll| false for the current route;
  \item the route node is marked \verb|.scalingUncleared|
    (\verb|navierSupercriticalScalingNode|,
    \verb|navierAveragedEquationAllObstructionsNode|), and the roadmap forces the
    averaged-equation obstruction to be examined before any replacement route
    (\verb|currentNavierRoadmap_orders_averaged_obstruction_before_replacement|).
\end{itemize}
So nothing here is a regularity theorem; what is proved is a precise inventory of
\emph{exactly which obligations are unmet}. That inventory is the route's real
asset.

\subsection{Possible avenues still open}

\paragraph{N1. Supply a critical-norm control witness.} The most direct missing
witness is a proof that the $L^3$ (or $\dot H^{1/2}$) critical norm is
controlled from finite energy plus the dissipation budget
(\verb|criticalInitialDataFamily| $+$ \verb|criticalNormControlWitness|). This is
the single witness that flips \verb|criticalNorm| from \verb|.blocked| toward
\verb|.cleared|.

\paragraph{N2. A scale-invariant upgrade witness.} Independently, supply a
\verb|scaleInvariantUpgradeWitness| against the
\verb|supercriticalConcentrationFamily| --- a quantitative statement that
concentration at small scales cannot outrun dissipation. This is the analytic
heart and the place the supercritical wall is felt most sharply.

\paragraph{N3. Represent the averaged-equation obligations honestly.} Each of the
four obligations maps to a required witness
(\verb|NavierAveragedEquationObligation.requiredWorldInterface|). The avenue is
to represent them one at a time --- averaged momentum, Reynolds-stress closure,
energy-inequality transport, critical-scaling transport --- so that
\verb|represented| grows from $[]$ toward the full list. Each representation is a
finite, checkable target rather than a leap to regularity.

\subsection{Conjectured workarounds for the supercritical wall}

\begin{conjecture}[Speculative --- Averaging as a critical-norm gain]
\label{conj:ns-averaging-gain}
The averaging operator, applied to the divergence-free velocity, gains exactly
the half-derivative needed to move the controlled diagnostic from the
supercritical $L^2$ regime to the critical $L^3$ / $\dot H^{1/2}$ regime, with a
constant that is uniform under the parabolic rescaling. Formally, this would
supply both \verb|criticalNormControlWitness| and
\verb|scaleInvariantUpgradeWitness| from a single averaging estimate.
\end{conjecture}

\noindent\emph{Why it might hold.} Averaging is a smoothing operation; if the
Reynolds-stress closure can be written so that the averaged nonlinearity is
strictly subcritical, the gate's two blocked layers could clear together.
\emph{Why it might fail.} The supercritical scaling is scale-invariant: a
smoothing that helps at one scale generically loses exactly as much at the dual
scale. Tao's averaged-Navier--Stokes blowup constructions are the cautionary
tale --- a sufficiently general averaged equation can blow up, so the conjecture
is responsible only if the specific averaging here is shown to respect a sign or
structure that the blowup constructions violate.

\begin{conjecture}[Speculative --- Reynolds closure with a coercive remainder]
\label{conj:ns-reynolds}
The Reynolds-stress closure (\verb|reynoldsStressClosureWitness|) can be chosen
so that its remainder term is \emph{coercive} in the critical norm: the
averaged-equation energy identity gains a strictly positive critical-norm
dissipation absent from the unaveraged identity. This coercivity, if uniform
under rescaling, is precisely the
\verb|criticalScalingCompatibility| obligation.
\end{conjecture}

\noindent\emph{Why it might hold.} Closures are not unique; a closure tuned for
coercivity rather than physical fidelity is a legitimate proof device.
\emph{Why it might fail.} Coercivity in the critical norm is essentially the
problem itself; the conjecture risks assuming the conclusion unless the
coercive remainder is derived from a genuinely new estimate.

\subsection{Where to point creative effort next}

The single highest-leverage target is
\textbf{the scale-invariant upgrade witness (avenue N2), pursued via
Conjecture~\ref{conj:ns-averaging-gain}}. Reason: the Lean gate makes plain that
the route's verdict hinges on exactly one analytic fact --- whether anything in
the averaged formulation beats supercritical concentration uniformly under
rescaling. Every other obligation is downstream bookkeeping that becomes
tractable once that estimate exists, and is provably blocked while it does not.
The recommended push: write the candidate averaging estimate as a Lean
\verb|scaleInvariantUpgradeWitness| against a concrete
\verb|supercriticalConcentrationFamily|, and stress-test it against a discretized
analogue of Tao's averaged blowup --- if the estimate survives the family that is
designed to make averaged Navier--Stokes blow up, it is a genuine candidate; if
it fails on that family, the failure pinpoints which structural sign the real
Navier--Stokes nonlinearity must supply that the abstract averaged equation
lacks. Either outcome converts a vague ``supercriticality'' into a specific,
checkable obstacle.

\subsection*{PLN lens}

\begin{table}[h]
\centering
\small
\begin{tabular}{lcccc}
\toprule
Central node & STV $\langle s,c\rangle$ & ITV $[L,U]$ & Progress & Status \\
\midrule
Scale-invariant upgrade ($L^2\!\to\!L^3$ uniform) & $\langle 0.05, 0.18\rangle$ & $[0.009,\,0.829]$ & --- & open crux (supercritical wall) \\
Averaged-equation obligations represented & $\langle 0.10, 0.25\rangle$ & $[0.025,\,0.775]$ & --- & all missing \\
Promotion-gate / obligation audit (the asset) & $\langle 0.97, 0.85\rangle$ & $[0.8245,\,0.9745]$ & --- & built \\
\midrule
\textbf{Route (through)} & $\langle 0.04, 0.21\rangle$ & $[0.0084,\,0.7984]$ & $\sim$30\% & \\
\bottomrule
\end{tabular}
\caption{PLN summary for the averaged-equation Navier--Stokes route. The audit
layer is high-strength (it is fully built and machine-checked); the route's own
expectation is very low because it is provably scaling-uncleared and the decisive
estimate confronts the supercritical-scaling wall.}
\end{table}

\noindent Through-route: $\langle 0.04, 0.21\rangle$, ITV
$[0.04\cdot 0.21,\;0.04\cdot 0.21+(1-0.21)] = [0.0084,\,0.7984]$,
PROGRESS $\sim 30\%$. The low strength reflects the proved scaling-uncleared
status: as it stands the route does not prove regularity, and the supercritical
wall is exactly the place every known energy-method route also stalls. Progress
is lower than 4CP/PNP because the analytic obligations are not merely unproved
but largely unrepresented --- the route has mapped \emph{that} they are missing,
not yet \emph{how} to supply them.
